% floodfusion -- IGARSS 2027 full paper (4 pages), DRAFT.  Build: see paper/igarss27/README.md.
% Every number is taken from a committed report under reports/; its source is named in a comment beside it.
% "% TODO(r3)" marks the ladder's rung 3, filled from reports/v13_ladder.json once the queue has finished.
% The official IGARSS 2027 template comes with the submission portal's paper kit; this uses IEEEtran (conference).
\documentclass[conference]{IEEEtran}
\usepackage[T1]{fontenc}
\usepackage{newtxtext,newtxmath}
\usepackage{graphicx}
\usepackage{booktabs}
\usepackage{xcolor}
\usepackage{url}
\usepackage[hidelinks]{hyperref}   % the shared references.bib uses \href
\newcommand{\tstar}{\ensuremath{\tau^{*}}}
\newcommand{\todo}[1]{\textcolor{red}{[#1]}}

\begin{document}

\title{How Close Is Sentinel-1 Flood Segmentation to Its Reference?\\
Labels, Threshold and the Self-Agreement of Copernicus EMS Maps}

\author{\IEEEauthorblockN{Teo Clerici Jurado}
\IEEEauthorblockA{BSc in AI \& Data Science, H-FARM College / University of Chichester\\
Code, data cards and reports: \url{https://github.com/teo-clerk/floodfusion}}}

\maketitle

\begin{abstract}
Deep flood segmentation from Sentinel-1 is usually reported on benchmarks whose labels were drawn on optical images,
while the reference that emergency services use is the Copernicus Emergency Management Service (CEMS) rapid mapping.
We score models zero-shot against the CEMS delineations of 38 flood activations of 2023--26 drawn on radar
satellites, and ask what limits the agreement. Training labels move the score most: the same network reaches 0.397
macro IoU when trained on radar-drawn flood labels and 0.258 on labels drawn mostly on optical images. Our best model
reaches 0.417 at its cost-derived threshold and 0.468 at each flood's best threshold in hindsight, while a
cross-fitted recalibration gains only 0.002. We then measure the reference itself: two CEMS delineations of the same
area at most 24\,h apart agree at 0.520 macro IoU over 34 activations, and per activation the model's score tracks
that self-agreement (Spearman $\rho = 0.51$, exploratory). Fine-tuning on CEMS's own labels, selected by a
preregistered rule, \todo{TODO(r3)}. The frozen model is under a preregistered, sealed blind test and maps new
floods automatically.
\end{abstract}

\begin{IEEEkeywords}
Flood mapping, Sentinel-1, SAR, Copernicus EMS, deep learning, evaluation, reference data.
\end{IEEEkeywords}

\section{Introduction}
Flood maps from synthetic aperture radar are the first product after a flood: Sentinel-1 sees through cloud and at
night. Deep segmentation models are usually trained and evaluated on Sen1Floods11~\cite{bonafilia2020sen1floods11},
whose hand labels were drawn mostly on Sentinel-2. An analyst in an emergency, however, compares against the CEMS
rapid-mapping delineations~\cite{cems2026rapidmapping}, drawn by photo-interpreters on the images available, often
radar. How well do models agree with that reference, and what limits the agreement?

We separate three candidate limits: the \emph{training labels}, the \emph{decision threshold}, and the
\emph{reference itself}. The contributions are:
(i) a zero-shot benchmark on 38 radar-drawn CEMS activations of 2023--26;
(ii) to our knowledge, the first measurement of how well two CEMS delineations of the same flood agree, and how
model scores relate to it;
(iii) an adaptation ladder, recalibration then fine-tuning on CEMS labels, with an adoption rule fixed before any
result; and
(iv) an open model frozen for a preregistered, sealed blind test on activations no model had seen, with a live
service that maps new floods from the frozen weights.

\section{Data}
\textbf{Inputs.} Sentinel-1 RTC backscatter (VV, VH) from Microsoft Planetary Computer~\cite{planetarycomputer2022}
for the post-event pass and a same-orbit pre-event reference; a seasonal reference, the median backscatter of the
same orbit around the same date a year earlier; Copernicus DEM terrain; ESA WorldCover~\cite{zanaga2022worldcover}
flags; and Sentinel-2 L2A where it is clear.
% src: reports/protocol_v1_2.md; MODEL_CARD.md ("v1.2's model"); DATA_CARD.md

\textbf{Training labels.} Kuro Siwo~\cite{bountos2024kurosiwo} flood labels, drawn on radar, plus the GEOID-Flood
events~\cite{chiriaco2026geoid} before the benchmark that Kuro Siwo lacks: 7,006 chips from 138 floods in 81 groups
of floods within 100\,km. Sen1Floods11 serves only as a comparison of label sources.
% src: reports/v1_2_deviations.md (pooled training set row); MODEL_CARD.md

\textbf{Benchmark.} Every CEMS flood activation EMSR656--898 (2023--26) whose delineation can be paired with a
Sentinel-1 pass within 24\,h, by a rule committed before scoring. The primary stratum holds the products drawn on
radar satellites: 38 activations with reference flood in their scored region (EMSR664, our case study, excluded).
Pixels outside the CEMS area of interest, in radar shadow or layover are not scored.
% src: reports/protocol_v1_1.md §2.3; reports/v12_cems_rebuilt.md (K = 38)

\section{Methods}
\textbf{Model.} A U-Net whose SAR and terrain encoder is fused with an optical encoder by per-pixel cross-attention,
with the seasonal reference and an availability flag as extra SAR channels; windows without usable optical use a
separate calibration. Logits are temperature-scaled on held-out location groups. The decision threshold
$\tstar = 0.25$ follows from a stated cost ratio: a missed flood pixel costs three false alarms.
% src: MODEL_CARD.md (v1.2's model); reports/protocol_v1_2.md (tau*, cost ratio 3)

\textbf{Metric.} IoU of the flood class per activation, pooled over its tiles, averaged over activations (macro).

\textbf{The reference ceiling.} For every area of interest of the seen activations with two consecutive delivered
delineations whose source images are at most 24\,h apart (the benchmark's own pairing window), both are rasterised
on a 10\,m grid over the region both observed, and their IoU is computed. Pairs are averaged within activations,
then across. A sensitivity row drops pairs in which the later product repeats more than 5\,\% of the earlier flood
unchanged. Floods also change in 24\,h, so this is a lower bound on the reference's precision, not a noise estimate.
% src: reports/protocol_v1_3.md §2.3, §5; reports/v1_3_deviations.md (ceiling set, sensitivity row)

\textbf{The adaptation ladder.} Activations are grouped by 100\,km single linkage (34 groups) and split into five
folds. R1 is the zero-shot model at \tstar. R2 refits only the threshold, as the argmin of $3\,\mathrm{FN} +
\mathrm{FP}$ on the other folds. R3 fine-tunes the model on the other folds' CEMS chips (30,834 in all), with a
learning rate chosen on an inner split. A rung is adopted only if it gains at least 0.005 macro IoU over the one
below \emph{and} is better on more than half of the 38 activations. Both rules were committed before any rung was
computed.
% src: reports/protocol_v1_3.md §2.1, §3; reports/v1_3_deviations.md

\section{Results}
\subsection{Labels move the score most}
On the 38 activations, the same fusion network scores 0.397 trained on Kuro Siwo's radar-drawn flood labels, 0.338
on the water labels of the same chips, and 0.258 on Sen1Floods11, below a physics baseline of thresholds on
backscatter change (0.371). Across architectures trained on the same labels, scores differ by about 0.01.
% src: reports/cems_benchmark.md (rows: Fusion S1F11 0.258; Kuro Siwo flood 0.397; water 0.338; physics 0.371)
% src: reports/v12_cems_rebuilt.md (v1.2 0.417, fair TerraMind 0.405, v1.1 best 0.407)

\subsection{The threshold leaves about 0.05, and it cannot be learned}
Our model (v1.2) reaches 0.417 macro IoU at \tstar, against 0.405 for TerraMind-base~\cite{jakubik2025terramind}
fine-tuned on the same labels and 0.321 for the Copernicus Global Flood Monitoring ensemble~\cite{salamon2021gfm} on
the 34 activations it covers (0.419 for v1.2 there). At each activation's best threshold in hindsight it would reach
0.468. Yet R2, the cross-fitted cost-optimal threshold, gains only +0.002 and is better on 26 of 38 activations:
not adopted. The best threshold varies between floods, not between our training set and CEMS.
% src: reports/v12_cems_rebuilt.md (0.417, 0.405, 0.468; GFM 0.321 vs 0.419 on 34); reports/v12_cems_seen_tests.md
% src: reports/v13_ladder_r2.md (R1 0.421, R2 0.423, +0.002, 26 of 38; seed 42)

\subsection{Copernicus agrees with itself at 0.52}
Over 79 product pairs in 34 activations, two CEMS delineations of the same area at most 24\,h apart agree at 0.520
macro IoU (median pair 0.498); 0.515 between two radar-drawn products, and 0.509 in the sensitivity row. Models at
0.42 are therefore near the level at which the reference agrees with itself. Per activation (Fig.~\ref{fig:ceiling}),
the model scores below the reference's self-agreement in 20 of 27 activations, and the two are associated
(Spearman $\rho = 0.51$). This analysis was chosen after both were known; floods that are hard for an analyst may
simply be hard for a model too.
% src: reports/reference_ceiling.md (79 pairs, 34 activations, 0.520, 0.498, radar-radar 0.515, sensitivity 0.509)
% src: paper/igarss27/figures/ceiling_scatter.json (n 27, rho 0.509, p 0.0067, below 20)

\begin{figure}[t]
\centering
\includegraphics[width=\columnwidth]{figures/ceiling_scatter.pdf}
\caption{Per seen activation: the v1.2 model's IoU at \tstar{} against CEMS (y) and the IoU between two CEMS
delineations of the same area at most 24\,h apart (x). Dashed: equal agreement. Exploratory.}
\label{fig:ceiling}
\end{figure}

\subsection{Fine-tuning on Copernicus's own labels}
\todo{TODO(r3): R3's macro IoU over the 38 activations, its gain over the adopted rung, the number of activations it
improves, and the rule's verdict; the per-cluster table; whether a v1.3 model was frozen.}
% src: reports/v13_ladder.md, reports/v13_ladder.json

\section{The blind test and the live service}
The models were frozen (tag \texttt{v1.2.0-freeze}, checkpoints pinned by sha256) before any flood of the blind
set, every CEMS flood activation from EMSR899 on, was scored. Scores are sealed in a hash-chained ledger until 20
activations in 15 clusters are scored, or 31 December 2027. Three hypotheses, set against GFM, the fair TerraMind and
the previous model, form a Holm family; a later, separate hypothesis H8 enters a v1.3 model only if the ladder
adopts one. \todo{TODO(status): the ledger's counts at submission.}
% src: reports/protocol_v1_2.md (§2.2, sealing, stopping rule, Holm family); reports/protocol_v1_3.md §4

The frozen weights also run a live service. Every six hours a scheduled job maps each new Sentinel-1 pass over
flood areas from the activation's outline alone, never reading a CEMS product. It rebuilds the benchmark's inputs:
on a replayed benchmark tile all 14 input bands were identical, and a CPU run gave the stored GPU run's IoU. Blind
floods appear there, map only, after the v1.3 decision.
% src: reports/live_replay.md; reports/protocol_v1_3.md §6; reports/v1_3_deviations.md (2026-11-11)

\section{Conclusion}
Against Copernicus EMS, Sentinel-1 flood segmentation is limited first by its training labels, then by a threshold
that varies between floods and cannot be refit from other floods, and finally by a reference that agrees with
itself at about 0.52. Scores near 0.42 should be read against that ceiling, not against 1. \todo{TODO(r3): one
sentence on fine-tuning.} Code, reports, the frozen weights and the blind test's ledger are public.

\section*{Acknowledgment}
Contains modified Copernicus Sentinel data and Copernicus Emergency Management Service information (2023--26),
\textcopyright{} European Union. The code and text were developed with AI assistance (Claude Code); the design
decisions, verification and results were reviewed by the author.

\bibliographystyle{IEEEtran}
\bibliography{../references}

\end{document}
